\chapter{Vektor}\label{ch:g10:vectors}

Sebuah \omterm{def:g10:vectors:vector}{vektor} menyandikan perpindahan: sebuah arah dan sebuah panjang,
tanpa memedulikan titik awalnya. \omterm{def:g10:vectors:vector}{Vektor} memberi bukti yang bersih bagi
fakta geometris dan, begitu \omterm{def:g10:coordgeom:system}{koordinat} masuk ke dalam gambar, menyusutkannya
menjadi perhitungan kecil. \omterm{def:g10:vectors:vector}{Vektor} akan menjadi perkakas penting pada
\cref{ch:g10:lines} dan, kelak, bagi hasil kali skalar pada \cref{ch:g11:scal}.

\section{Translasi dan vektor}

\begin{definition}[Vektor]\label{def:g10:vectors:vector}
Diberikan dua titik $A$ dan $B$, \emph{vektor}\index{vektor} $\vect{AB}$
menyatakan perpindahan dari $A$ ke $B$: vektor itu mempunyai sebuah
\emph{arah} (arah garis $(AB)$, beserta arah perjalanan sepanjang
garis itu, dari $A$ menuju $B$) dan sebuah
\emph{panjang} (jarak $AB$, ditulis juga $\norm{\vect{AB}}$).
Dua vektor disebut \emph{sama} apabila keduanya menyandikan perpindahan yang
sama: $\vect{AB} = \vect{CD}$ berarti bahwa menerapkan kedua perpindahan itu
pada titik mana pun memberi hasil yang sama.
\end{definition}

\begin{proposition}[Vektor yang sama dan jajargenjang]\label{prop:g10:vectors:parallelogram}
$\vect{AB} = \vect{CD}$ tepat ketika $ABDC$ (menurut urutan itu!) adalah
jajargenjang, mungkin jajargenjang yang gepeng --- setara dengan itu, ketika
$[AD]$ dan $[BC]$ mempunyai \omterm{prop:g10:coordgeom:midpoint}{titik tengah} yang sama.
\end{proposition}
\admitted

\begin{omfigure}
\begin{tikzpicture}[scale=1.0]
  \coordinate (A) at (0,0);
  \coordinate (B) at (2.2,0.8);
  \coordinate (C) at (1.0,-1.2);
  \coordinate (D) at (3.2,-0.4);
  \draw[omExo, thin] (A) -- (C);
  \draw[omExo, thin] (B) -- (D);
  \draw[-Stealth, omDef, thick] (A) -- (B)
    node[midway, above, font=\small] {$\vect{AB}$};
  \draw[-Stealth, omDef, thick] (C) -- (D)
    node[midway, below, font=\small] {$\vect{CD}$};
  \fill (A) circle (1.5pt) node[left, font=\small] {$A$};
  \fill (B) circle (1.5pt) node[right, font=\small] {$B$};
  \fill (C) circle (1.5pt) node[left, font=\small] {$C$};
  \fill (D) circle (1.5pt) node[right, font=\small] {$D$};
\end{tikzpicture}

{\small \omterm{def:g10:vectors:vector}{Vektor} yang sama $\vect{AB} = \vect{CD}$: arah sama, jurusan sama,
panjang sama. Segi empat $ABDC$ adalah jajargenjang.}
\end{omfigure}

\begin{notation}
Sebuah \omterm{def:g10:vectors:vector}{vektor} sering dinamai dengan satu huruf, $\vec u$ atau $\vec v$, bila
titik ujungnya tidak penting. \emph{\omterm{def:g10:vectors:vector}{Vektor} nol} $\vec 0 = \vect{AA}$ adalah
perpindahan yang tidak menggerakkan apa pun. \emph{Lawan} dari
$\vect{AB}$ adalah $\vect{BA}$: panjangnya sama, jurusannya berlawanan. Kita tulis
$\vect{BA} = -\vect{AB}$.
\end{notation}

\section{Menjumlahkan vektor}

\begin{definition}[Jumlah dua vektor]\label{def:g10:vectors:sum}
Yang disebut \emph{jumlah}\index{vektor!jumlah} $\vec u + \vec v$ adalah perpindahan
yang diperoleh dengan melakukan $\vec u$ lalu $\vec v$.
\end{definition}

\begin{theorem}[Relasi Chasles]\label{thm:g10:vectors:chasles}
Untuk sebarang tiga titik $A$, $B$, $C$:
\[
\vect{AB} + \vect{BC} = \vect{AC}.
\]
\end{theorem}

\begin{proof}
Pergi dari $A$ ke $B$, lalu dari $B$ ke $C$, adalah perpindahan yang membawa
$A$ ke $C$: itulah perpindahan yang disandikan oleh $\vect{AC}$.
\end{proof}

\begin{remark}[Aturan jajargenjang]
Untuk menjumlahkan dua \omterm{def:g10:vectors:vector}{vektor} yang digambar dari titik yang sama,
$\vect{AB} + \vect{AC}$, lengkapkan jajargenjang $ABDC$: \omterm{def:g10:vectors:sum}{jumlahnya} adalah
diagonal $\vect{AD}$. Memang $\vect{AC} = \vect{BD}$, sehingga
$\vect{AB} + \vect{AC} = \vect{AB} + \vect{BD} = \vect{AD}$ menurut relasi
Chasles.
\end{remark}

\begin{omfigure}
\begin{tikzpicture}[scale=1.05]
  \coordinate (A) at (0,0);
  \coordinate (B) at (2.4,0.5);
  \coordinate (C) at (3.4,2.0);
  \draw[-Stealth, omDef, thick] (A) -- (B)
    node[midway, below, font=\small] {$\vec u$};
  \draw[-Stealth, omThm, thick] (B) -- (C)
    node[midway, right, font=\small] {$\vec v$};
  \draw[-Stealth, omMeth, thick] (A) -- (C)
    node[midway, above left, font=\small] {$\vec u + \vec v$};
\end{tikzpicture}\hspace{1.4cm}%
\begin{tikzpicture}[scale=1.05]
  \coordinate (A) at (0,0);
  \coordinate (B) at (2.4,0.5);
  \coordinate (C) at (1.0,1.5);
  \coordinate (D) at (3.4,2.0);
  \draw[omExo, thin] (B) -- (D) -- (C);
  \draw[-Stealth, omDef, thick] (A) -- (B)
    node[midway, below, font=\small] {$\vect{AB}$};
  \draw[-Stealth, omThm, thick] (A) -- (C)
    node[midway, left, font=\small] {$\vect{AC}$};
  \draw[-Stealth, omMeth, thick] (A) -- (D)
    node[pos=0.75, below right, font=\small] {$\vect{AD}$};
  \fill (D) circle (1.2pt) node[right, font=\small] {$D$};
\end{tikzpicture}

{\small Dua gambar penjumlahan: ujung ke pangkal (relasi Chasles, kiri) dan
aturan jajargenjang (kanan).}
\end{omfigure}

\begin{example}[Menyederhanakan dengan Chasles]\label{ex:g10:vectors:chasles}
Sederhanakan $\vect{MN} + \vect{NP} + \vect{PQ}$: merangkai perpindahannya
memberi $\vect{MQ}$. Sederhanakan $\vect{AB} - \vect{AC}$: tulis ulang
selisihnya sebagai \omterm{def:g10:vectors:sum}{jumlah} dengan \omterm{def:g10:vectors:vector}{vektor} lawannya,
\[
\vect{AB} - \vect{AC} = \vect{AB} + \vect{CA}
= \vect{CA} + \vect{AB} = \vect{CB}.
\]
\end{example}

\section{Mengalikan vektor dengan bilangan}

\begin{definition}[Kelipatan skalar]\label{def:g10:vectors:scalar}
Misalkan $\vec u$ \omterm{def:g10:vectors:vector}{vektor} taknol dan $k$ \omterm{def:g10:numbers:sets}{bilangan real}. \omterm{def:g10:vectors:vector}{Vektor}
$k\vec u$\index{vektor!kelipatan skalar} mempunyai arah yang sama dengan
$\vec u$, panjang $\abs{k} \times \norm{\vec u}$, dan berjurusan sama
dengan $\vec u$ jika $k > 0$, berlawanan jika $k < 0$. Untuk $k = 0$,
$0\vec u = \vec 0$.
\end{definition}

\begin{definition}[Kekolinearan]\label{def:g10:vectors:collinear}
Dua \omterm{def:g10:vectors:vector}{vektor} $\vec u$ dan $\vec v$ disebut
\emph{kolinear}\index{vektor kolinear} apabila keduanya mempunyai arah yang
sama, yaitu apabila $\vec v = k \vec u$ untuk suatu \omterm{def:g10:numbers:sets}{bilangan real} $k$ (atau
salah satunya $\vec 0$).
\end{definition}

\begin{remark}
Kekolinearan adalah bahasa \omterm{def:g10:vectors:vector}{vektor} untuk kesejajaran dan kesegarisan:
\begin{itemize}
  \item garis $(AB)$ dan $(CD)$ sejajar tepat ketika $\vect{AB}$
        dan $\vect{CD}$ \omterm{def:g10:vectors:collinear}{kolinear};
  \item titik $A$, $B$, $C$ segaris tepat ketika $\vect{AB}$ dan
        $\vect{AC}$ \omterm{def:g10:vectors:collinear}{kolinear}.
\end{itemize}
\end{remark}

\section{Koordinat vektor}

\begin{definition}[Koordinat sebuah vektor]\label{def:g10:vectors:coords}
Pada sebuah \omterm{def:g10:coordgeom:system}{sistem koordinat}, \emph{\omterm{def:g10:coordgeom:system}{koordinat}} \omterm{def:g10:vectors:vector}{vektor} $\vect{AB}$
adalah bilangan yang memerikan perpindahannya:
\[
\vect{AB}\,(x_B - x_A,\ y_B - y_A).
\]
Sebuah \omterm{def:g10:vectors:vector}{vektor} $\vec u\,(a, b)$ menggeser setiap titik sejauh $a$ satuan
mendatar dan $b$ satuan tegak.
\end{definition}

\begin{proposition}[Berhitung dengan koordinat]\label{prop:g10:vectors:coordrules}
Misalkan $\vec u\,(a, b)$ dan $\vec v\,(c, d)$, serta $k \in \R$.
\begin{enumerate}
  \item $\vec u = \vec v$ tepat ketika $a = c$ dan $b = d$;
  \item $\vec u + \vec v$ berkoordinat $(a + c,\ b + d)$;
  \item $k\vec u$ berkoordinat $(ka,\ kb)$;
  \item $\norm{\vec u} = \sqrt{a^2 + b^2}$ (pada \omterm{def:g10:coordgeom:system}{sistem ortonormal}).
\end{enumerate}
\end{proposition}

\begin{proof}
1--3 menyatakan ulang, \omterm{def:g10:coordgeom:system}{koordinat} demi \omterm{def:g10:coordgeom:system}{koordinat}, apa yang dilakukan
perpindahannya: misalnya melakukan $\vec u$ lalu $\vec v$ menggeser sebuah
titik secara mendatar sejauh $a$ lalu $c$, jadi $a + c$ seluruhnya. Butir 4
adalah rumus jarak \cref{thm:g10:coordgeom:distance} yang diterapkan pada
ruas garis yang mewakili $\vec u$.
\end{proof}

\begin{theorem}[Kriteria kekolinearan]\label{thm:g10:vectors:det}
Dua \omterm{def:g10:vectors:vector}{vektor} $\vec u\,(a, b)$ dan $\vec v\,(c, d)$ \omterm{def:g10:vectors:collinear}{kolinear} jika dan
hanya jika
\[
ad - bc = 0 .
\]
\end{theorem}

\begin{proof}
Jika $\vec v = k\vec u$, maka $c = ka$ dan $d = kb$, sehingga
$ad - bc = a(kb) - b(ka) = 0$. Hal yang sama berlaku jika $\vec u = k \vec v$
atau jika salah satu \omterm{def:g10:vectors:vector}{vektornya} nol.

Sebaliknya, andaikan $ad - bc = 0$ dengan $\vec u \neq \vec 0$, katakan
$a \neq 0$ (kasus $b \neq 0$ serupa). Ambil $k = \frac{c}{a}$; maka
$c = ka$, dan $ad = bc = b(ka)$ memberi, setelah dibagi $a \neq 0$,
$d = kb$. Jadi $\vec v = k\vec u$.
\end{proof}

\begin{example}\label{ex:g10:vectors:det}
Apakah $\vec u\,(3, -2)$ dan $\vec v\,(-6, 4)$ \omterm{def:g10:vectors:collinear}{kolinear}? Hitunglah
$ad - bc = 3 \times 4 - (-2)\times(-6) = 12 - 12 = 0$: ya, dan memang
$\vec v = -2\vec u$. Untuk $\vec u\,(3, -2)$ dan $\vec w\,(1, 5)$:
$3 \times 5 - (-2) \times 1 = 17 \neq 0$: tidak \omterm{def:g10:vectors:collinear}{kolinear}.
\end{example}

\begin{method}[Kesegarisan dan kesejajaran]\label{met:g10:vectors:alignment}
Untuk membuktikan bahwa tiga titik $A$, $B$, $C$ segaris:
\begin{enumerate}
  \item hitunglah \omterm{def:g10:coordgeom:system}{koordinat} $\vect{AB}$ dan $\vect{AC}$;
  \item periksalah kriteria $ad - bc = 0$ pada \cref{thm:g10:vectors:det};
  \item simpulkan: kedua \omterm{def:g10:vectors:vector}{vektornya} \omterm{def:g10:vectors:collinear}{kolinear} dan berbagi titik $A$, jadi
        ketiga titiknya segaris.
\end{enumerate}
Perhitungan yang sama dengan $\vect{AB}$ dan $\vect{CD}$ membuktikan bahwa
garis $(AB)$ dan $(CD)$ sejajar.
\end{method}

\begin{example}\label{ex:g10:vectors:alignment}
Misalkan $A(1, 2)$, $B(3, 5)$ dan $C(7, 11)$. Maka $\vect{AB}\,(2, 3)$ dan
$\vect{AC}\,(6, 9)$, serta $2 \times 9 - 3 \times 6 = 18 - 18 = 0$: jadi
ketiga titiknya segaris (bahkan $\vect{AC} = 3\vect{AB}$).
\end{example}

\begin{omfigure}
\begin{tikzpicture}[scale=0.62]
  \draw[thin, gray!40] (-0.4,-0.4) grid (8.4,5.4);
  \draw[-Stealth] (-0.6,0) -- (8.6,0);
  \draw[-Stealth] (0,-0.6) -- (0,5.6);
  \coordinate (A) at (1,1);
  \coordinate (B) at (3,2);
  \coordinate (C) at (7,4);
  \draw[-Stealth, omDef, very thick] (A) -- (B)
    node[midway, above left=-2pt, font=\small] {$\vect{AB}$};
  \draw[-Stealth, omThm, thick] (A) -- (C)
    node[pos=0.8, below right=-2pt, font=\small] {$\vect{AC} = 3\vect{AB}$};
  \fill (A) circle (2pt) node[above left, font=\small] {$A$};
  \fill (B) circle (2pt) node[above left, font=\small] {$B$};
  \fill (C) circle (2pt) node[above left, font=\small] {$C$};
\end{tikzpicture}

{\small Titik yang segaris: $\vect{AC}$ adalah kelipatan skalar $\vect{AB}$.}
\end{omfigure}

\section{Latihan}

\begin{exercise}[$\star$]\label{exo:g10:vectors:1}
Sederhanakan dengan relasi Chasles:
\[
\vect{AB} + \vect{BD}, \qquad
\vect{KL} + \vect{LM} + \vect{MK}, \qquad
\vect{AC} - \vect{BC}, \qquad
\vect{AB} + \vect{CA}.
\]
\end{exercise}

\begin{exercise}[$\star$]\label{exo:g10:vectors:2}
Misalkan $A(2, 1)$, $B(5, 3)$, $C(-1, 4)$. Hitunglah \omterm{def:g10:coordgeom:system}{koordinat}
$\vect{AB}$, $\vect{BC}$, $\vect{AC}$, lalu periksa pada \omterm{def:g10:coordgeom:system}{koordinatnya} bahwa
$\vect{AB} + \vect{BC} = \vect{AC}$.
\end{exercise}

\begin{exercise}[$\star$]\label{exo:g10:vectors:3}
Misalkan $\vec u\,(2, -3)$ dan $\vec v\,(-1, 4)$. Berikan \omterm{def:g10:coordgeom:system}{koordinat}
$\vec u + \vec v$, $3\vec u$, $2\vec u - \vec v$, lalu hitunglah
$\norm{\vec u}$.
\end{exercise}

\begin{exercise}[$\star$]\label{exo:g10:vectors:4}
Pada setiap hal berikut, sebutkan apakah $\vec u$ dan $\vec v$ \omterm{def:g10:vectors:collinear}{kolinear}:
\[
\text{(a) } \vec u\,(4, 6),\ \vec v\,(6, 9);
\qquad
\text{(b) } \vec u\,(2, -5),\ \vec v\,(-4, 10);
\qquad
\text{(c) } \vec u\,(3, 1),\ \vec v\,(1, 3).
\]
\end{exercise}

\begin{exercise}[$\star$]\label{exo:g10:vectors:5}
Misalkan $A(0, 2)$, $B(4, 0)$, $C(6, 3)$. Carilah \omterm{def:g10:coordgeom:system}{koordinat} titik
$D$ sehingga $ABCD$ menjadi jajargenjang, yaitu sehingga
$\vect{AB} = \vect{DC}$.
\end{exercise}

\begin{exercise}[$\star\star$]\label{exo:g10:vectors:6}
Misalkan $A(-2, 1)$, $B(1, 3)$, $C(4, -1)$.
\begin{enumerate}
  \item Hitunglah \omterm{def:g10:coordgeom:system}{koordinat} titik $M$ sehingga
        $\vect{AM} = \vect{AB} + \vect{AC}$.
  \item Segi empat apakah $ABMC$ itu? Berikan alasannya.
\end{enumerate}
\end{exercise}

\begin{exercise}[$\star\star$]\label{exo:g10:vectors:7}
Misalkan $A(1, -1)$, $B(4, 1)$, $C(10, 5)$. Apakah titik $A$, $B$, $C$
segaris? Pertanyaan yang sama untuk $A(0, 3)$, $B(2, 2)$, $C(8, -1)$.
\end{exercise}

\begin{exercise}[$\star\star$]\label{exo:g10:vectors:8}
Misalkan $A(1, 2)$, $B(5, 4)$, $C(6, 1)$, $D(2, -1)$. Tunjukkan bahwa $(AB)$
dan $(CD)$ sejajar, lalu bahwa $ABCD$ adalah jajargenjang. Pemeriksaan yang
mana mengakibatkan yang lain?
\end{exercise}

\begin{exercise}[$\star\star$]\label{exo:g10:vectors:9}
Misalkan $ABC$ sebuah segitiga, dan misalkan $I$ titik yang ditentukan oleh
$\vect{AI} = \frac23 \vect{AB}$. Dengan $A(0,0)$, $B(6, 3)$, $C(2, 5)$
(agar perhitungannya tetap sederhana):
\begin{enumerate}
  \item hitunglah \omterm{def:g10:coordgeom:system}{koordinat} $I$;
  \item misalkan $J$ titik dengan $\vect{CJ} = \frac23\vect{CB}$; hitunglah
        \omterm{def:g10:coordgeom:system}{koordinat} $J$;
  \item tunjukkan bahwa $(IJ)$ sejajar dengan $(AC)$.
\end{enumerate}
\end{exercise}

\begin{exercise}[$\star\star\star$]\label{exo:g10:vectors:10}
Misalkan $ABCD$ sebarang segi empat, dan misalkan $I$, $J$, $K$, $L$ \omterm{prop:g10:coordgeom:midpoint}{titik
tengah} $[AB]$, $[BC]$, $[CD]$, $[DA]$. Dengan memakai \omterm{def:g10:coordgeom:system}{koordinat}
$A(x_A, y_A)$, dan seterusnya, tunjukkan bahwa $\vect{IJ} = \vect{LK}$, lalu
simpulkan bahwa \omterm{prop:g10:coordgeom:midpoint}{titik tengah} sisi pada \emph{sebarang} segi empat selalu
membentuk jajargenjang.
\end{exercise}

\section{Soal: Aljabar anak panah}

\begin{problem}[{Soal akhir pekan --- senam Chasles, titik berat
yang dibuktikan ulang dalam dua baris, dan gaya dalam keseimbangan:
untuk apa sebenarnya vektor itu}]
\label{pb:g10:vectors:1}
Di jilid sebelumnya sudah dibuktikan bahwa ketiga garis berat sebuah
segitiga bertemu pada dua pertiga panjangnya --- dengan sebuah jajargenjang
cerdik yang tersembunyi di dalam segitiganya (soal akhir pekan tentang titik
berat di jilid sebelumnya). \omterm{def:g10:vectors:vector}{Vektor} membuktikannya ulang dalam
dua baris, dan menghadiahkan perpotongan di satu titik itu secara
cuma-cuma. Untuk itulah aljabar baru ini ada: teorema menjadi perhitungan dengan
anak panah. Soal ini melatih perhitungannya (terutama Chasles),
menyajikan bukti dua barisnya, memperluas Varignon
(\cref{exo:g10:vectors:10}), dan berakhir di tempat \omterm{def:g10:vectors:vector}{vektor} dilahirkan:
gaya dan kecepatan.

\medskip
\noindent\textbf{Bagian I --- Senam Chasles.}
\begin{enumerate}
  \item Sederhanakan dengan relasi Chasles
        (\cref{thm:g10:vectors:chasles}):
        $\vect{AB} + \vect{BC} + \vect{CD}$; \quad
        $\vect{MA} - \vect{MB}$; \quad
        $\vect{AB} + \vect{CD} + \vect{BC} + \vect{DA}$.
  \item Kiat titik asal: tunjukkan bahwa untuk \emph{sebarang} titik $O$,
        berlaku $\vect{AB} = \vect{OB} - \vect{OA}$.
  \item Buktikan kedua ciri \omterm{prop:g10:coordgeom:midpoint}{titik tengah} $I$ dari
        $[AB]$:
        \[
        \vect{IA} + \vect{IB} = \vec 0
        \qquad\text{dan}\qquad
        \vect{OI} = \tfrac12\left(\vect{OA} + \vect{OB}\right)
        \ \text{untuk sebarang } O .
        \]
  \item Dengan memakai $\vect{AB} = \vect{DC}$, temukan kembali titik sudut
        keempat $D$ jajargenjang $ABCD$ dengan $A(-1, 2)$,
        $B(3, 4)$, $C(6, 0)$ --- lalu bandingkan dengan
        jawaban lewat kiat \omterm{prop:g10:coordgeom:midpoint}{titik tengah} pada \cref{pb:g10:coordgeom:1}.
  \item Menambahkan sebuah \omterm{def:g10:vectors:vector}{vektor} tetap $\vec u$ pada setiap titik pada
        bidang menjalankan transformasi yang mana --- transformasi
        yang ditemukan dengan dua cermin pada soal akhir pekan tentang dua
        cermin di jilid sebelumnya dan dengan dua setengah putaran pada soal
        akhir pekan tentang setengah putaran di jilid sebelumnya? Hitunglah
        \omterm{def:g10:functions:function}{peta} $(1, 2)$ oleh $\vec u = (3, -1)$.
\end{enumerate}

\medskip
\noindent\textbf{Bagian II --- Titik berat, bukti ketiga.}
Definisikan titik $G$ sebuah segitiga $ABC$ oleh syarat yang
anggun berikut
\[
\vect{GA} + \vect{GB} + \vect{GC} = \vec 0 .
\]
\begin{enumerate}[resume]
  \item Dengan kiat titik asal, tunjukkan bahwa syarat itu
        menentukan $G$ secara tunggal:
        $\vect{OG} = \frac13\left(\vect{OA} + \vect{OB} +
        \vect{OC}\right)$ untuk sebarang $O$ --- sehingga \omterm{def:g10:coordgeom:system}{koordinat}
        $G$ adalah \emph{rata-rata} \omterm{def:g10:coordgeom:system}{koordinat} ketiga titik sudutnya
        (soal akhir pekan tentang titik berat di jilid sebelumnya sudah
        melihatnya secara numeris).
  \item Buktikan $\vect{AG} = \frac13\left(\vect{AB} +
        \vect{AC}\right)$, dan, dengan menulis $K$ untuk \omterm{prop:g10:coordgeom:midpoint}{titik tengah}
        $[BC]$, buktikan $\vect{AK} = \frac12\left(\vect{AB} +
        \vect{AC}\right)$. Simpulkan:
        $\vect{AG} = \frac23\,\vect{AK}$ --- yaitu teorema dua
        pertiga, terbukti ulang.
  \item Mengapa garis berat dari $B$ dan dari $C$ melalui
        $G$ yang \emph{sama} tanpa perhitungan tambahan?
        Bandingkan ketiga bukti teorema ini yang tersedia ---
        jajargenjang tersembunyi di jilid sebelumnya, \omterm{def:g10:coordgeom:system}{koordinat}, dan \omterm{def:g10:vectors:vector}{vektor}
        --- masing-masing dalam satu kalimat.
  \item Pemeriksaan numeris: $A(1, 1)$, $B(5, 2)$, $C(3, 6)$.
        Hitunglah $G$, lalu $K$, lalu periksalah
        $\vect{AG} = \frac23 \vect{AK}$ dengan kriteria
        kekolinearan (\cref{thm:g10:vectors:det}).
  \item Fisika membaca $G$ sebagai titik seimbang tiga massa
        yang sama di ketiga titik sudutnya. Lipatduakan massa di $A$ (massanya
        menjadi $2, 1, 1$): titik seimbangnya menjadi
        $\vect{OG'} = \frac{2\vect{OA} + \vect{OB} +
        \vect{OC}}{4}$. Hitunglah $G'$ untuk segitiga pada
        pertanyaan 9 lalu bandingkan letaknya dengan $G$.
\end{enumerate}

\medskip
\noindent\textbf{Bagian III --- Kekolinearan sedang bekerja.}
\begin{enumerate}[resume]
  \item Apakah $\vec u(3, -2)$ dan $\vec v(-6, 4)$ \omterm{def:g10:vectors:collinear}{kolinear}? Lalu
        $\vec w(2, 5)$ dan $\vec z(4, 9)$? Putuskan dengan
        kriteria determinannya.
  \item Apakah titik $A(1, 2)$, $B(3, 5)$, $C(7, 11)$
        segaris?
  \item Melampaui Varignon (\cref{exo:g10:vectors:10}): pada sebarang
        segi empat $ABCD$, \emph{bimedian} menghubungkan
        \omterm{prop:g10:coordgeom:midpoint}{titik tengah} sisi yang berhadapan ($[IK]$ dan $[JL]$). Dengan
        memakai \omterm{def:g10:vectors:vector}{vektor} posisi ($\vect{OI} = \frac12(\vect{OA} +
        \vect{OB})$, dan seterusnya), tunjukkan bahwa kedua bimedian itu
        mempunyai \omterm{prop:g10:coordgeom:midpoint}{titik tengah} yang sama, dengan \omterm{def:g10:vectors:vector}{vektor} posisi
        $\frac14\left(\vect{OA} + \vect{OB} + \vect{OC} +
        \vect{OD}\right)$: yaitu titik berat segi empat itu.
  \item Thales dalam satu baris: $M$ dan $N$ adalah titik yang ditentukan
        oleh $\vect{AM} = \frac13\vect{AB}$ dan
        $\vect{AN} = \frac13\vect{AC}$. Hitunglah
        $\vect{MN}$ dalam $\vect{BC}$ lalu simpulkan
        kesejajaran serta perbandingannya.
  \item Perumumlah pertanyaan 14 ke perbandingan sebarang $k$, lalu
        nyatakan dalam satu kalimat bagaimana kalkulus \omterm{def:g10:vectors:vector}{vektor} itu
        menyerap teorema \omterm{prop:g10:coordgeom:midpoint}{titik tengah} dan
        teorema Thales sekaligus.
\end{enumerate}

\medskip
\noindent\textbf{Bagian IV --- Gaya dan kecepatan.}
\begin{enumerate}[resume]
  \item Sebuah perahu didorong arus $3$ km/jam ke timur
        sementara didayung $4$ km/jam ke utara. Berikan \omterm{def:g10:vectors:vector}{vektor}
        kecepatan resultannya, besarnya, lalu uraikan haluan yang
        sesungguhnya. (Sebuah tripel yang sangat tua muncul di sini.)
  \item Tiga gaya bekerja pada sebuah cincin: $\vec u(2, 1)$,
        $\vec v(-3, 2)$, $\vec w(1, -3)$. Hitunglah resultannya.
        Apa yang dilakukan cincin itu? Dan mengapa \omterm{def:g10:vectors:sum}{jumlah} anak panah
        sepanjang segi banyak \emph{tertutup} mana pun selalu $\vec 0$
        (Chasles)?
  \item Dua gaya $\vec u(4, -1)$ dan $\vec v(-1, 3)$ bekerja pada sebuah
        kait. Gaya ketiga apa yang menahan kait itu tetap seimbang?
  \item Seorang perenang menyeberangi sungai dengan menuju lurus ke seberang
        pada $2$ m/s; arusnya mengalir $1.5$ m/s. Berikan
        kecepatan resultan dan besarnya. Sungai itu selebar
        $50$ m: berapa lama penyeberangannya, dan seberapa
        jauh ke hilir perenang itu mendarat?
  \item Penutup: satu objek, tiga busana --- anak panah
        perpindahan, pasangan \omterm{def:g10:coordgeom:system}{koordinat}, gaya atau kecepatan. Daftarkan
        ketiga teorema klasik yang dibuktikan ulang soal ini dengan
        aljabar \omterm{def:g10:vectors:vector}{vektor}, lalu sebutkan satu besaran geometris yang belum
        dapat diukur oleh \omterm{def:g10:vectors:vector}{vektor} bab ini --- perkakas yang membereskannya
        tiba bersama hasil kali skalar, pada tahun berikutnya.

\end{enumerate}
\end{problem}
